% Annex A — perímetres, àrees i volums (horitzontal)
\renewcommand{\arraystretch}{2.0}\small
\newcommand{\aT}[2]{\begin{tabular}{@{}p{4.1cm}@{\hspace{0.5em}}c@{\hspace{0.5em}}c@{}}\toprule\textbf{#1}&\textbf{Perímetre}&\textbf{Àrea}\\\midrule #2 \\\bottomrule\end{tabular}}
\begin{minipage}[t]{0.49\linewidth}\centering
\aT{Figura plana}{
Quadrat (costat $l$) & $P=4l$ & $A=l^2$ \\
Rectangle (base $b$, alçada $h$) & $P=2b+2h$ & $A=b\,h$ \\
Paral·lelogram (base $b$, alçada $h$, costat $c$) & $P=2b+2c$ & $A=b\,h$ \\
Rombe (diagonals $D,d$; costat $l$) & $P=4l$ & $A=\D\frac{D\,d}{2}$ \\
Triangle (base $b$, alçada $h$) & $P=a+b+c$ & $A=\D\frac{b\,h}{2}$ \\
Triangle equilàter (costat $l$) & $P=3l$ & $A=\D\frac{\sqrt3}{4}\,l^2$ \\
Trapezi (bases $B,b$; costats $c,d$) & $P=B+b+c+d$ & $A=\D\frac{B+b}{2}\,h$}
\end{minipage}\hfill
\begin{minipage}[t]{0.49\linewidth}\centering
\aT{Figura plana}{
Polígon regular ($n$ costats, apotema $a$) & $P=n\,l$ & $A=\D\frac{P\,a}{2}$ \\
Cercle (radi $r$) & $L=2\pi r$ & $A=\pi r^2$ \\
Sector circular (radi $r$, angle $\theta$ en rad) & $L=\theta\,r$ & $A=\D\frac{\theta\,r^2}{2}$ \\
Corona circular (radis $R>r$) & $2\pi(R+r)$ & $A=\pi(R^2-r^2)$ \\
El·lipse (semieixos $a,b$) & $P\approx 2\pi\sqrt{\D\frac{a^2+b^2}{2}}$ & $A=\pi\,a\,b$ \\
Segment circular (radi $r$, angle $\theta$) & $L=\theta\,r$ & $A=\D\frac{r^2}{2}(\theta-\sin\theta)$ \\
Hexàgon regular (costat $l$) & $P=6l$ & $A=\D\frac{3\sqrt3}{2}\,l^2$}
\end{minipage}
\par\vspace{1.4em}
\renewcommand{\arraystretch}{2.05}
\newcommand{\vT}[2]{\begin{tabular}{@{}p{4.1cm}@{\hspace{0.5em}}c@{\hspace{0.5em}}c@{}}\toprule\textbf{#1}&\textbf{Àrea}&\textbf{Volum}\\\midrule #2 \\\bottomrule\end{tabular}}
\begin{minipage}[t]{0.49\linewidth}\centering
\vT{Cos geomètric}{
Cub (aresta $a$) & $A_T=6a^2$ & $V=a^3$ \\
Ortoedre (arestes $a,b,c$) & $A_T=2(ab+bc+ac)$ & $V=abc$ \\
Prisma (àrea base $A_B$, perímetre $P_B$) & $A_L=P_B\,h;\ A_T=2A_B+A_L$ & $V=A_B\,h$ \\
Piràmide (apotema $a_p$) & $A_L=\D\frac{P_B\,a_p}{2};\ A_T=A_B+A_L$ & $V=\D\frac{A_B\,h}{3}$ \\
Cilindre (radi $r$, alçada $h$) & $A_L=2\pi r h;\ A_T=2\pi r(h+r)$ & $V=\pi r^2h$}
\end{minipage}\hfill
\begin{minipage}[t]{0.49\linewidth}\centering
\vT{Cos geomètric}{
Con (radi $r$, generatriu $g$) & $A_L=\pi r g;\ A_T=\pi r(g+r)$ & $V=\D\frac{\pi r^2h}{3}$ \\
Tronc de con (radis $R,r$) & $A_L=\pi g(R+r)$ & $V=\D\frac{\pi h}{3}(R^2+Rr+r^2)$ \\
Esfera (radi $r$) & $A=4\pi r^2$ & $V=\D\frac{4\pi r^3}{3}$ \\
Semiesfera (radi $r$) & $A_T=3\pi r^2$ & $V=\D\frac{2\pi r^3}{3}$ \\
Calota esfèrica (alçada $h$) & $A=2\pi r h$ & $V=\D\frac{\pi h^2}{3}(3r-h)$}
\end{minipage}
\par\vspace{1.6em}
\renewcommand{\arraystretch}{1.5}\normalsize
\begin{minipage}[t]{0.32\linewidth}
\textbf{Teorema de Pitàgores}\par (triangle rectangle de catets $a,b$ i hipotenusa $c$):
\[c^2=a^2+b^2\]
\[c=\sqrt{a^2+b^2},\qquad a=\sqrt{c^2-b^2}\]
\end{minipage}\hfill
\begin{minipage}[t]{0.32\linewidth}
\textbf{Raons trigonomètriques}\par (triangle rectangle, angle $\alpha$):
\[\sin\alpha=\frac{\text{oposat}}{\text{hipotenusa}},\ \ \cos\alpha=\frac{\text{adjacent}}{\text{hipotenusa}}\]
\[\tan\alpha=\frac{\text{oposat}}{\text{adjacent}},\qquad \sin^2\alpha+\cos^2\alpha=1\]
\end{minipage}\hfill
\begin{minipage}[t]{0.32\linewidth}
\textbf{Fórmula de Heró}\par (semiperímetre $s=\frac{a+b+c}{2}$):
\[A=\sqrt{s(s-a)(s-b)(s-c)}\]
\textbf{Àrea d'un triangle amb dos costats:}
\[A=\tfrac12\,a\,b\,\sin\gamma\]
\end{minipage}
\normalsize
